{{دیوان
| عنوان = تقدیر (اِبلیس و یَزداں)
| ماخذ = ویکی ماخذ
| ماخذ_لنک = https://ur.wikisource.org/wiki/%D8%AA%D9%82%D8%AF%DB%8C%D8%B1_%28%D8%A7%D9%90%D8%A8%D9%84%DB%8C%D8%B3_%D9%88_%DB%8C%D9%8E%D8%B2%D8%AF%D8%A7%DA%BA%29
| مقابلہ = اقبال سائبر لائبریری (اقبال اکادمی پاکستان)
| مقابلہ_لنک = https://www.iqbalcyberlibrary.net
}}
<poem>
اِبلیس

اے خدائے کُن فکاں! مجھ کو نہ تھا آدم سے بَیر

آہ! وہ زندانیِ نزدیک و دُور و دیر و زُود

حرفِ ’اِستکبار‘ تیرے سامنے ممکن نہ تھا

ہاں، مگر تیری مشیّت میں نہ تھا میرا سجود

یزداں
کب کھُلا تجھ پر یہ راز، انکار سے پہلے کہ بعد؟

اِبلیس
بعد، اے تیری تجلّی سے کمالاتِ وجود!

یزداں
(فرشتوں کی طرف دیکھ کر)

پستیِ فطرت نے سِکھلائی ہے یہ حُجتّ اسے
کہتا ہے ’تیری مشیّت میں نہ تھا میرا سجود،

دے رہا ہے اپنی آزادی کو مجبوری کا نام
ظالم اپنے شُعلۂ سوزاں کو خود کہتا ہے دُود!

(ماخوذ از محی الدّین ابنِ عربیؒ)
</poem>
